\documentclass[11pt]{article}
\usepackage{amsmath,amssymb,amsfonts}
\usepackage{geometry}
\geometry{margin=1in}
\usepackage{graphicx}
\usepackage{booktabs}

\newcommand{\bxi}{\boldsymbol{\xi}}
\newcommand{\bx}{\mathbf{x}}
\newcommand{\avg}[1]{\langle #1 \rangle}

\title{Centroid Escape and the Structure of the Retrieval Boundary in LSR Dense Associative Memory}
\author{}
\date{}

\begin{document}
\maketitle

\section{Per-trial overlap distributions near the retrieval boundary}

\subsection{Setup and observables}

We study the LSR (Log-Sum-ReLU) energy on the $N$-sphere $|\bx|^2 = N$:
\begin{equation}
E(\bx) = -\frac{N}{b}\log\left(\sum_{\mu=1}^{P}\max\!\Big(0,\;1 - b + b\,\frac{\bxi^\mu \cdot \bx}{N}\Big)\right),
\label{eq:lsr}
\end{equation}
with $b = 2+\sqrt{2}$.  The hard wall is at $\varphi_c = (b-1)/b \approx 0.707$: pattern~$\mu$ is ``in support'' of state~$\bx$ when $\varphi_\mu \equiv \bxi^\mu\!\cdot\!\bx/N > \varphi_c$.

For each grid point $(\alpha, T)$, the v8m simulation generates $n_{\mathrm{dis}}$ independent disorder samples (random pattern sets $\{\bxi^\mu\}$).  Each sample produces two replicas (a, b) initialized near the target $\bxi^1$ with $\varphi_{\mathrm{init}}\in[0.75,\,1.0]$, equilibrated for $2^{14}$ Metropolis steps, then sampled for $2^{12}$ steps.  The time-averaged overlap $\varphi_a$, $\varphi_b$ and the maximum spurious overlap $\varphi_{\max}^{\mathrm{other}}$ are recorded per sample.

\subsection{Three populations in the CDF}

Figure~\ref{fig:cdf} shows the empirical CDF of the per-trial overlap $\varphi$ at selected $(\alpha, T)$ points.  At $T = 0.325$ ($\varphi_{\mathrm{eq}} = 0.90$), the distribution is \emph{not} a single Gaussian.  Instead, three discrete populations emerge:

\begin{center}
\begin{tabular}{lcccl}
\toprule
Population & $\varphi$ & $\sigma(\varphi)$ & $\varphi_{\max}^{\mathrm{other}}$ & Interpretation \\
\midrule
Retrieval & 0.901 & 0.003 & 0.54--0.66 & single-pattern equilibrium \\
Centroid  & 0.831 & 0.010 & 0.77--0.82 & two-pattern equilibrium \\
Escaped   & $< 0.75$ & broad & varies & multi-pattern or paramagnetic \\
\bottomrule
\end{tabular}
\end{center}

The centroid peak at $\varphi \approx 0.831$ is remarkably narrow ($\sigma \approx 0.01$) and \emph{constant across all $\alpha$ values}, even deep inside the retrieval phase ($\alpha = 0.200$).  Its defining signature is $\varphi_{\max}^{\mathrm{other}} \approx \varphi$: the system has nearly equal overlap with the target and one spurious pattern.  As $\alpha$ increases toward $\alpha_{\mathrm{th}}$, the weight of the centroid population grows from $\sim$7\% to $\sim$60\%, while the retrieval peak shrinks correspondingly.

\begin{figure}[t]
\centering
% \includegraphics[width=\textwidth]{cdf_phi.pdf}
\caption{Empirical CDF of the per-trial time-averaged overlap~$\varphi$ from v8m data.  Green: R-region points; orange: M-region; blue: P-region.  The dashed gray line marks $\varphi_c \approx 0.707$.  The two-step structure in the green and orange curves reflects the retrieval peak ($\varphi \approx 0.90$) and the centroid peak ($\varphi \approx 0.83$).}
\label{fig:cdf}
\end{figure}

\section{The centroid escape mechanism}

\subsection{Support overlap between patterns}

Consider the target pattern $\bxi^1$ and a spurious pattern~$\bxi^\mu$ with inter-pattern overlap $q_{\!1\mu} = \bxi^1\!\cdot\bxi^\mu/N$.  For random patterns on $S^{N-1}(\!\sqrt{N})$, $q_{\!1\mu}$ is approximately Gaussian with mean~0 and standard deviation $1/\sqrt{N}$.

At $N \approx 47$ and $P \in [20\,000,\,200\,000]$, the extreme-value statistics of the maximum overlap is
\begin{equation}
q_{\max} \approx \sqrt{\frac{2\ln P}{N}} \;\approx\; 0.66\text{--}0.77\,.
\end{equation}
The number of patterns with $q_{1\mu} > q^*$ is
$
n(q^*) = (P-1)\,\Phi^c\!\big(q^*\sqrt{N}\big),
$
where $\Phi^c$ is the complementary normal CDF.  For $q^* = 0.38$ (the value needed to produce the observed centroid at $\varphi = 0.831$; see below):
\begin{center}
\begin{tabular}{cccr}
\toprule
$\alpha$ & $N$ & $P$ & $n(q>0.38)$ \\
\midrule
0.200 & 50 & 20\,000 & 70 \\
0.220 & 49 & 43\,000 & 165 \\
0.240 & 47 & 88\,000 & 396 \\
0.265 & 46 & 201\,000 & 981 \\
\bottomrule
\end{tabular}
\end{center}

Thus, even at $\alpha = 0.200$---well below $\alpha_{\mathrm{th}} \approx 0.245$---tens to hundreds of spurious patterns have substantial overlap with the target.

\subsection{Geometry of the zero-barrier escape}

We work in the two-dimensional subspace spanned by the orthonormalized pair
\[
\mathbf{e}_1 = \frac{\bxi^1}{\sqrt{N}}, \qquad
\mathbf{e}_2 = \frac{\bxi^\mu/\sqrt{N} - q\,\mathbf{e}_1}{\sqrt{1-q^2}}\,.
\]
Projecting the state $\bx/\sqrt{N} = a\,\mathbf{e}_1 + b\,\mathbf{e}_2 + \mathbf{c}_\perp$, with $a^2 + b^2 + |\mathbf{c}_\perp|^2 = 1$, the overlaps are
\begin{equation}
\varphi_1 = a, \qquad \varphi_\mu = a\,q + b\sqrt{1-q^2}\,.
\label{eq:overlaps}
\end{equation}

The support conditions require $\varphi_1 > \varphi_c$ and $\varphi_\mu > \varphi_c$.  These define two half-planes in the $(a,b)$ disk (Fig.~\ref{fig:geometry}).  The overlap region---where both patterns are simultaneously in support---is non-empty whenever $q > 2\varphi_c^2 - 1 = 0$ (i.e., always for positive $q$).

\begin{figure}[t]
\centering
% \includegraphics[width=0.6\textwidth]{centroid_geometry.pdf}
\caption{Support geometry projected onto the $(\bxi^1, \bxi^\mu)$ plane for $q = 0.38$.  The unit circle represents the sphere constraint $a^2 + b^2 \leq 1$.  Blue shading: support of $\bxi^1$ ($a > \varphi_c$).  Red shading: support of $\bxi^\mu$ ($\varphi_\mu > \varphi_c$).  The escape path has two steps: \textcircled{1}~rotate~$\bx_\perp$ at fixed $\varphi_1$, bringing $\bxi^\mu$ into support with zero energy cost; \textcircled{2}~slide downhill to the centroid.}
\label{fig:geometry}
\end{figure}

The escape path from the retrieval state to the centroid has two stages (Fig.~\ref{fig:geometry}):

\paragraph{Stage~1: perpendicular rotation ($\Delta E = 0$).}
At the retrieval state, $a = \varphi_{\mathrm{eq}} \approx 0.90$ and $b \approx 0$.  The overlap with $\bxi^\mu$ is $\varphi_\mu \approx \varphi_{\mathrm{eq}}\,q \approx 0.34$, far below $\varphi_c$---only $\bxi^1$ is in support.

The system can increase~$b$ (rotating $\bx_\perp$ toward~$\bxi^\mu_\perp$) while keeping $a = \varphi_1$ fixed.  This moves the state vertically in the $(a,b)$ plane without changing $\bxi^1$'s ReLU contribution, so \emph{the energy is exactly constant}.  When $b$ reaches
\begin{equation}
b_{\mathrm{entry}} = \frac{\varphi_c - \varphi_{\mathrm{eq}}\,q}{\sqrt{1-q^2}}
\;\approx\; \frac{0.707 - 0.90 \times 0.38}{0.925} \;=\; 0.394,
\end{equation}
the spurious pattern enters support ($\varphi_\mu = \varphi_c$).  At this point, $a^2 + b^2 = 0.81 + 0.16 = 0.97 < 1$, confirming that the state remains on the sphere.

\paragraph{Stage~2: downhill slide to the centroid ($\Delta E < 0$).}
Once both patterns are in support, the sum inside the logarithm of Eq.~\eqref{eq:lsr} increases, and the energy \emph{decreases}.  The system slides to the two-pattern centroid at
\begin{equation}
\varphi_{\mathrm{cen}} = \sqrt{\frac{1+q}{2}} \;\approx\; \sqrt{\frac{1.38}{2}} \;=\; 0.831\,,
\label{eq:phi_cen}
\end{equation}
in quantitative agreement with the observed mid-peak (Table above).

\subsection{Energetics}

The energy per neuron at the retrieval state (one pattern in support) and at the centroid (two patterns, equal overlap) is:
\begin{align}
\frac{E_{\mathrm{ret}}}{N} &= -\frac{1}{b}\ln\!\big(1 - b + b\,\varphi_{\mathrm{eq}}\big)
= -\frac{1}{b}\ln(0.659) = 0.122\,, \\
\frac{E_{\mathrm{cen}}}{N} &= -\frac{1}{b}\ln\!\big(2(1 - b + b\,\varphi_{\mathrm{cen}})\big)
= -\frac{1}{b}\ln(0.846) = 0.049\,.
\end{align}

The centroid has \emph{lower} energy by $\Delta E/N = -0.073$.  Including the entropy on the sphere, $S/N \approx \tfrac{1}{2}\ln(1-\varphi^2)$, the free energy difference is
\begin{equation}
\frac{\Delta F}{N} = \frac{F_{\mathrm{cen}} - F_{\mathrm{ret}}}{N} \approx -0.15
\qquad (T = 0.325)\,,
\end{equation}
confirming that the centroid is \emph{thermodynamically preferred}.

\subsection{The barrier is entropic, not energetic}

The energy barrier along the escape path is exactly zero: Stage~1 has $\Delta E = 0$ and Stage~2 has $\Delta E < 0$.  The system can reach the centroid without ever increasing its energy.

What prevents instantaneous escape is an \emph{entropic barrier}: the rotation in Stage~1 must align the $(N\!-\!1)$-dimensional perpendicular component $\bx_\perp$ with a \emph{specific} direction ($\bxi^\mu_\perp$) in a high-dimensional space.  The Metropolis random walk explores $O(1)$ new directions per step with step size $\sigma = 2.4\,T/\!\sqrt{N}$, so the typical time to discover the escape direction scales as $O(N/\sigma^2)$.

For $N = 47$, $T = 0.325$: $N/\sigma^2 \approx 3\,400$ steps.  With $2^{14} + 2^{12} \approx 20\,000$ Metropolis steps per trial, the system has $\sim$6 independent attempts to find each escape direction.  With $n(q > 0.38) \sim 70$--$1\,000$ competing patterns, the total escape probability is substantial, explaining the observed centroid fractions of 7--60\%.

The escape probability depends on \emph{both} the number of competing patterns (which grows with $P$ and hence $\alpha$) and the allocated MC time, making the centroid fraction a kinetic quantity that would change with longer equilibration.

\subsection{Relation to the classical Hopfield model}

In the classical Hopfield model (Amit, Gutfreund \& Sompolinsky 1987), the analog of the centroid state is the ``mixture state''---a metastable configuration with equal overlap $m_n$ with an odd number~$n$ of patterns.  However, there are fundamental differences:

\begin{enumerate}
\item \textbf{Origin.}  Classical mixture states arise from the nonlinear coupling in the Hebb rule and exist as saddle-point solutions of the mean-field equations.  LSR centroid states arise geometrically from the \emph{compact support} of the Epanechnikov kernel: they form wherever support regions of distinct patterns overlap on the $N$-sphere.

\item \textbf{Energy.}  Classical mixture states have energy \emph{higher} than retrieval states (Amit et al., \S4, point~4).  LSR centroids have energy \emph{lower} than the single-pattern retrieval minimum, because the log-sum gains from having two active ReLU terms.

\item \textbf{Barrier.}  In the classical model, retrieval and mixture states are separated by barriers of order~$N$.  In LSR, the barrier is \emph{zero in energy}---purely entropic.

\item \textbf{Role of $P$.}  Classical mixture states exist for finite~$p$; their number grows combinatorially but they remain metastable.  LSR centroids are the \emph{global minimum} for disorder realizations with nearby patterns, and their prevalence grows with $P$ through extreme-value statistics.
\end{enumerate}

This mechanism is unique to associative memories with compact-support kernels (LSR / Epanechnikov) and has no analog in the LSE (Gaussian) model, where the exponential tail provides global support and no geometric overlap regions exist.

\subsection{Implications for the phase diagram}

The centroid mechanism explains the three anomalies observed in the v8 phase diagram near the curved retrieval boundary at $\alpha > \alpha_{\mathrm{th}}$:

\begin{enumerate}
\item \textbf{Retrieval phase extends beyond theory.}  The theoretical boundary $\alpha_c(T)$ assumes the \emph{only} minimum is the single-pattern retrieval state.  In practice, the average $\avg{\varphi}$ is pulled up by the centroid population (with $\varphi \approx 0.83$), creating apparent retrieval beyond the boundary.

\item \textbf{$q_{\mathrm{EA}} > \varphi^2$.}  At the centroid, both replicas are attracted to the same two-pattern minimum: $\bx_a \cdot \bx_b / N$ gains a contribution from the correlated perpendicular components (both aligned with $\bxi^\mu$), giving $q > \varphi_1^2$.

\item \textbf{$\varphi - \varphi_{\max}^{\mathrm{other}} \approx 0.1$ at the boundary.}  At the centroid, the target and spurious overlaps are nearly equal ($\varphi \approx \varphi_{\max}^{\mathrm{other}} \approx 0.83$), so their difference is small.  The disorder average, mixing centroid and retrieval populations, gives $\varphi - \varphi_{\max}^{\mathrm{other}} \sim 0.1$.
\end{enumerate}

\end{document}
